% ============================================================
\setpartpreamble{\partintro
  {Custom SQL expressions, aggregates, run-time metadata, scaffolding and a second provider:
  the tools for very complex plugins.}
  {
    \item Chapter~\ref{ch:custom-sql-expression-nodes}: Custom SQL expression nodes
    \item Chapter~\ref{ch:aggregate-translators}: Aggregate translators
    \item Chapter~\ref{ch:property-level-fluent-api-and}: Property-level Fluent API and
    runtime annotations
    \item Chapter~\ref{ch:scaffolding-in-practice}: Scaffolding in practice
    \item Chapter~\ref{ch:what-belongs-in-the-core}: What belongs in the core
    \item Chapter~\ref{ch:a-second-provider-sql-server}: A second provider: SQL Server
  }}
\part{Deep extensions}\label{part:IX}

% ------------------------------------------------------------
\chapter{Custom SQL expression nodes}\label{ch:custom-sql-expression-nodes}

Chapter~\ref{ch:translating-methods-and-members} builds SQL from existing building blocks:
functions, constants, casts. That is enough for most plugins. A \textbf{custom expression
node} (a subclass of \code{Sql\allowbreak{}Expression}) becomes necessary when at least one
of the following applies:

\begin{datatable}{H{0.80}H{1.20}}
  \textbf{Reason} & \textbf{Example} \\
  \tablehead
  The syntax is not a function call & Operators such as \code{@@}, \code{<\allowbreak{}->},
  \code{IS DISTINCT FROM}, keywords such as \code{AT TIME ZONE} \\
  \tablerule
  The nullability semantics differ & \code{IS DISTINCT FROM} is never NULL, even when an
  operand is NULL \\
  \tablerule
  Other parts should recognize the node & Optimizations that only apply to this expression
  \\
\end{datatable}

The sample implements the null-safe comparison
\code{EF.\allowbreak{}Functions.\allowbreak{}Is\allowbreak{}Not\allowbreak{}Distinct\allowbreak{}From(\allowbreak{}a,\allowbreak{}
b)} and \code{Is\allowbreak{}Distinct\allowbreak{}From(\allowbreak{}a,\allowbreak{} b)}.
Npgsql does not offer it as a function; it meets all three reasons.

\section{Anatomy of a node}

\begin{datatable}{H{0.78}H{1.26}H{0.96}}
  \textbf{Member} & \textbf{Purpose} & \textbf{Consequence if missing} \\
  \tablehead
  \code{Visit\allowbreak{}Children} / \code{Update} & Visitors can replace children;
  unchanged children yield the same instance & EF throws at the first visitor \\
  \tablerule
  \code{Equals} / \code{Get\allowbreak{}Hash\allowbreak{}Code} & Structural comparison for
  the query cache and optimizations & Cache misses, wrong deduplication \\
  \tablerule
  \code{Print} & Readable representation in DebugView and error messages & Unclear errors \\
  \tablerule
  \code{Quote} & C\# code for precompiled queries and NativeAOT & No AOT support \\
  \tablerule
  SQL generator & Writes the node as SQL & \textbf{Silently} missing SQL \\
  \tablerule
  Nullability processor & Knows the NULL semantics & Exception during translation \\
\end{datatable}

The last two rows also apply to combinations with other plugins: if a second plugin replaces
the generator or the processor, the node's SQL is silently missing or the query throws.
\code{Composition\allowbreak{}Matrix\allowbreak{}Tests} shows both
(chapter~\ref{ch:registering-services}).

\begin{sourcecode}[Query/\allowbreak{}Expressions/\allowbreak{}Distinct\allowbreak{}From\allowbreak{}Expression.\allowbreak{}cs (\allowbreak{}shortened)]{csharp}
public class DistinctFromExpression : SqlExpression
{
    private static ConstructorInfo? _quotingConstructor;

    public DistinctFromExpression(SqlExpression left, SqlExpression right, bool negated,
        RelationalTypeMapping? typeMapping)
        : base(typeof(bool), typeMapping)
    {
        // ... ThrowIfNull ...
        Left = left;
        Right = right;
        IsNegated = negated;
    }

    public virtual SqlExpression Left { get; }
    public virtual SqlExpression Right { get; }
    public virtual bool IsNegated { get; }

    protected override Expression VisitChildren(ExpressionVisitor visitor)
        => Update((SqlExpression)visitor.Visit(Left), (SqlExpression)visitor.Visit(Right));

    public virtual DistinctFromExpression Update(SqlExpression left, SqlExpression right)
        => left == Left && right == Right
            ? this
            : new DistinctFromExpression(left, right, IsNegated, TypeMapping);

    public override Expression Quote()
        => New(
            _quotingConstructor ??= typeof(DistinctFromExpression).GetConstructor(
                [typeof(SqlExpression), typeof(SqlExpression), typeof(bool),
                 typeof(RelationalTypeMapping)])!,
            Left.Quote(), Right.Quote(), Constant(IsNegated),
            RelationalExpressionQuotingUtilities.QuoteTypeMapping(TypeMapping));

    protected override void Print(ExpressionPrinter expressionPrinter)
    {
        expressionPrinter.Visit(Left);
        expressionPrinter.Append(IsNegated ? " IS DISTINCT FROM " : " IS NOT DISTINCT FROM ");
        expressionPrinter.Visit(Right);
    }

    public override bool Equals(object? obj)
        => obj is DistinctFromExpression other
           && (ReferenceEquals(this, other)
               || base.Equals(other) && Left.Equals(other.Left)
                  && Right.Equals(other.Right) && IsNegated == other.IsNegated);

    public override int GetHashCode()
        => HashCode.Combine(base.GetHashCode(), Left, Right, IsNegated);
}
\end{sourcecode}

The translator creates the node and aligns the type mappings of both sides: if one side has
a mapping (a column, say), the other gets the same one. That way a constant becomes the
matching literal:

\begin{sourcecode}[Query/\allowbreak{}Distinct\allowbreak{}From\allowbreak{}Translator.\allowbreak{}cs (\allowbreak{}excerpt)]{csharp}
var definition = method.GetGenericMethodDefinition();
if (definition != IsNotDistinctFromMethod && definition != IsDistinctFromMethod)
    return null;

var left = arguments[1];                        // arguments[0] is DbFunctions
var right = arguments[2];
var typeMapping = ExpressionExtensions.InferTypeMapping(left, right)
                  ?? _typeMappingSource.FindMapping(left.Type);

return new DistinctFromExpression(
    _sql.ApplyTypeMapping(left, typeMapping)!,
    _sql.ApplyTypeMapping(right, typeMapping)!,
    negated: definition == IsDistinctFromMethod,
    _typeMappingSource.FindMapping(typeof(bool)));
\end{sourcecode}

\section{Extending the SQL generator}

The \code{Query\allowbreak{}Sql\allowbreak{}Generator} writes SQL using a visitor. It hands
unknown nodes to
\code{Expression\allowbreak{}Visitor.\allowbreak{}Visit\allowbreak{}Extension}, which only
visits the children. Without an override there is no error, but \textbf{SQL without the
operator}. The plugin derives from the Npgsql generator:

\begin{sourcecode}[Query/\allowbreak{}Internal/\allowbreak{}Tenancy\allowbreak{}Npgsql\allowbreak{}Query\allowbreak{}Sql\allowbreak{}Generator.\allowbreak{}cs (\allowbreak{}shortened)]{csharp}
public class TenancyNpgsqlQuerySqlGenerator : NpgsqlQuerySqlGenerator
{
    public TenancyNpgsqlQuerySqlGenerator(QuerySqlGeneratorDependencies dependencies,
        IRelationalTypeMappingSource typeMappingSource,
        bool reverseNullOrderingEnabled, Version postgresVersion)
        : base(dependencies, typeMappingSource, reverseNullOrderingEnabled, postgresVersion)
    {
    }

    protected override Expression VisitExtension(Expression extensionExpression)
        => extensionExpression switch
        {
            DistinctFromExpression e => VisitDistinctFrom(e),
            _ => base.VisitExtension(extensionExpression),
        };

    protected virtual Expression VisitDistinctFrom(DistinctFromExpression expression)
    {
        VisitOperand(expression.Left);
        Sql.Append(expression.IsNegated ? " IS DISTINCT FROM " : " IS NOT DISTINCT FROM ");
        VisitOperand(expression.Right);
        return expression;
    }
}
\end{sourcecode}

The generator is created through
\code{IQuery\allowbreak{}Sql\allowbreak{}Generator\allowbreak{}Factory}, a singleton single
service. The plugin replaces it with its own factory taking the same dependencies as
Npgsql's (\code{Query\allowbreak{}Sql\allowbreak{}Generator\allowbreak{}Dependencies}, type
mapping source, \code{INpgsql\allowbreak{}Singleton\allowbreak{}Options}).

\section{Extending the nullability processor}

After translation, EF optimizes NULL comparisons. To do so it must know for every node
whether it can become NULL. Unknown nodes end up in
\code{Visit\allowbreak{}Custom\allowbreak{}Sql\allowbreak{}Expression}, whose base
implementation throws an exception. The processor hangs off a chain of three classes, all of
which must be replaced:

\begin{sourcecode}[Query/\allowbreak{}Internal/\allowbreak{}Tenancy\allowbreak{}Npgsql\allowbreak{}Sql\allowbreak{}Nullability\allowbreak{}Processor.\allowbreak{}cs (\allowbreak{}shortened)]{csharp}
public class TenancyNpgsqlSqlNullabilityProcessor : NpgsqlSqlNullabilityProcessor
{
    // Constructor: (…ProcessorDependencies, …ProcessorParameters) : base(...)

    protected override SqlExpression VisitCustomSqlExpression(
        SqlExpression sqlExpression, bool allowOptimizedExpansion, out bool nullable)
    {
        if (sqlExpression is DistinctFromExpression distinctFrom)
        {
            // Visit children (they can be optimized), their nullability does not matter
            var left = Visit(distinctFrom.Left, out _)!;
            var right = Visit(distinctFrom.Right, out _)!;

            nullable = false;                   // IS [NOT] DISTINCT FROM is never NULL
            return distinctFrom.Update(left, right);
        }

        return base.VisitCustomSqlExpression(
            sqlExpression, allowOptimizedExpansion, out nullable);
    }
}

public class TenancyNpgsqlParameterBasedSqlProcessor : NpgsqlParameterBasedSqlProcessor
{
    protected override Expression ProcessSqlNullability(
        Expression selectExpression, ParametersCacheDecorator parametersDecorator)
        => new TenancyNpgsqlSqlNullabilityProcessor(Dependencies, Parameters)
            .Process(selectExpression, parametersDecorator);
}

// + TenancyNpgsqlParameterBasedSqlProcessorFactory
//   : IRelationalParameterBasedSqlProcessorFactory
\end{sourcecode}

\begin{sourcecode}[Infrastructure/\allowbreak{}Tenancy\allowbreak{}Options\allowbreak{}Extension.\allowbreak{}cs › Apply\allowbreak{}Services (\allowbreak{}addition)]{csharp}
// 4) Deep extension: custom SQL expression node
services.Replace(ServiceDescriptor
    .Singleton<IQuerySqlGeneratorFactory, TenancyNpgsqlQuerySqlGeneratorFactory>());
services.Replace(ServiceDescriptor
    .Scoped<IRelationalParameterBasedSqlProcessorFactory,
        TenancyNpgsqlParameterBasedSqlProcessorFactory>());
\end{sourcecode}

\begin{important}[The price of custom nodes]
  Custom nodes require replacing two single services and deriving from internal Npgsql
  classes. Two plugins that both introduce custom nodes can only be combined if one derives
  from the other. Hence: function composition first; custom nodes only when one of the three
  reasons applies. And the replacement belongs prominently in the documentation.
\end{important}

\begin{note}[Experimental API]
  \code{Quote} and
  \code{Relational\allowbreak{}Expression\allowbreak{}Quoting\allowbreak{}Utilities} belong
  to the experimental API for precompiled queries. The compiler reports their use as EF9100;
  the sample suppresses it project-wide and documents this in
  \code{Directory.\allowbreak{}Build.\allowbreak{}props}.
\end{note}

\begin{sourcecode}[tests/\allowbreak{}.\allowbreak{}.\allowbreak{}.\allowbreak{}/\allowbreak{}Deep\allowbreak{}Query\allowbreak{}Tests.\allowbreak{}cs (\allowbreak{}excerpt)]{csharp}
[Fact]
public void IsDistinctFrom_in_projection_passes_nullability_processor()
{
    // Without a custom nullability processor, EF would throw here
    using var ctx = Contexts.Create();

    var sql = ctx.Customers
        .Select(c => EF.Functions.IsDistinctFrom(c.Name, (string?)null))
        .ToQueryString();

    Assert.Contains("IS DISTINCT FROM", sql);
}
\end{sourcecode}

% ------------------------------------------------------------
\chapter{Aggregate translators}\label{ch:aggregate-translators}

Aggregates such as \code{Sum} or \code{string\_\allowbreak{}agg} combine several rows,
typically after \code{Group\allowbreak{}By}. An
\code{IAggregate\allowbreak{}Method\allowbreak{}Call\allowbreak{}Translator} receives,
instead of a single argument, an \code{Enumerable\allowbreak{}Expression} that describes
\emph{what} is being aggregated:

\begin{datatable}{H{0.70}H{1.09}H{1.21}}
  \textbf{Property} & \textbf{Meaning} & \textbf{Origin in LINQ} \\
  \tablehead
  \code{Selector} & The aggregated expression & \code{g.\allowbreak{}Select(\allowbreak{}o
  => o.\allowbreak{}Total > 0)} \\
  \tablerule
  \code{Is\allowbreak{}Distinct} & Distinct values only &
  \code{g.\allowbreak{}Select(\allowbreak{}.\allowbreak{}.\allowbreak{}.\allowbreak{}).\allowbreak{}Distinct(\allowbreak{})}
  \\
  \tablerule
  \code{Predicate} & Filter before the aggregate & \code{g.\allowbreak{}Where(\allowbreak{}o
  => o.\allowbreak{}Id > 10)} \\
  \tablerule
  \code{Orderings} & Order (for order-dependent aggregates) &
  \code{g.\allowbreak{}Order\allowbreak{}By(\allowbreak{}.\allowbreak{}.\allowbreak{}.\allowbreak{})}
  \\
\end{datatable}

The sample adds \code{EF.\allowbreak{}Functions.\allowbreak{}Every(\allowbreak{}values)} for
PostgreSQL's \code{bool\_\allowbreak{}and}:

\begin{sourcecode}[Query/\allowbreak{}Every\allowbreak{}Aggregate\allowbreak{}Translator.\allowbreak{}cs (\allowbreak{}excerpt)]{csharp}
public virtual SqlExpression? Translate(MethodInfo method, EnumerableExpression source,
    IReadOnlyList<SqlExpression> arguments,
    IDiagnosticsLogger<DbLoggerCategory.Query> logger)
{
    if (method != EveryMethod || source.Selector is not SqlExpression selector)
        return null;

    // Npgsql's factory implements DISTINCT, FILTER (WHERE ...) and ORDER BY in aggregates
    if (_sql is not NpgsqlSqlExpressionFactory npgsql)
        return null;

    var boolean = _typeMappingSource.FindMapping(typeof(bool))!;

    // bool_and returns NULL for an empty set, Enumerable.All returns true
    return _sql.Coalesce(
        npgsql.AggregateFunction("bool_and", [selector], source,
            nullable: true, argumentsPropagateNullability: [false],
            typeof(bool), boolean),
        _sql.Constant(true, boolean),
        boolean);
}
\end{sourcecode}

\begin{sourcecode}[Generated SQL for g.\allowbreak{}Where(\allowbreak{}o => o.\allowbreak{}Id > 10).\allowbreak{}Select(\allowbreak{}o => o.\allowbreak{}Total > 0)]{sql}
SELECT COALESCE(bool_and(o."Total" > 0.0) FILTER (WHERE o."Id" > 10), TRUE)
FROM "Orders" AS o
GROUP BY o."TenantId"
\end{sourcecode}

\begin{important}[Semantics before syntax]
  An aggregate translator must match the semantics of the .NET side, not just the database
  syntax. Typical pitfalls: empty sets (NULL instead of 0, true or an empty string), NULL
  values in the input, and predicates. Npgsql itself aligns \code{string.\allowbreak{}Join}
  with \code{string\_\allowbreak{}agg} using \code{COALESCE}; the sample follows this
  pattern.
\end{important}

Databases without a \code{FILTER} clause (such as SQL Server) emulate predicates with
\code{CASE WHEN predicate THEN value END}. A provider-neutral aggregate translator must know
both; that is why this one lives in the PostgreSQL package.

% ------------------------------------------------------------
\chapter{Property-level Fluent API and runtime annotations}\label{ch:property-level-fluent-api-and}

\section{Property builder}

So far the plugin assumes that the tenant property is called \code{Tenant\allowbreak{}Id}
and that \code{ITenant\allowbreak{}Owned} is implemented. With a Fluent API method on
\code{Property\allowbreak{}Builder} it works without either, for example with existing
models using a property \code{Mandant\allowbreak{}Id} or an explicit interface
implementation:

\begin{sourcecode}[Usage]{csharp}
modelBuilder.Entity<Contract>().Property(c => c.Mandant).IsTenantProperty();
\end{sourcecode}

\begin{sourcecode}[Metadata/\allowbreak{}Tenancy\allowbreak{}Property\allowbreak{}Builder\allowbreak{}Extensions.\allowbreak{}cs (\allowbreak{}core package)]{csharp}
public static PropertyBuilder IsTenantProperty(
    this PropertyBuilder builder, bool isTenant = true)
{
    ArgumentNullException.ThrowIfNull(builder);

    if (isTenant)
        builder.Metadata.SetAnnotation(TenancyAnnotationNames.IsTenantProperty, true);
    else
        builder.Metadata.RemoveAnnotation(TenancyAnnotationNames.IsTenantProperty);

    return builder;
}

public static PropertyBuilder<TProperty> IsTenantProperty<TProperty>(
    this PropertyBuilder<TProperty> builder, bool isTenant = true)
    => (PropertyBuilder<TProperty>)IsTenantProperty((PropertyBuilder)builder, isTenant);

public static IConventionPropertyBuilder? IsTenantProperty(
    this IConventionPropertyBuilder builder,
    bool isTenant = true,
    bool fromDataAnnotation = false)
    => builder.HasAnnotation(TenancyAnnotationNames.IsTenantProperty,
           isTenant ? true : null, fromDataAnnotation) is null
        ? null
        : builder;
\end{sourcecode}

The method lives in the core package because it only sets metadata. The conventions of both
provider packages find the tenant property through the same core function:

\begin{sourcecode}[Metadata/\allowbreak{}Tenancy\allowbreak{}Entity\allowbreak{}Type\allowbreak{}Extensions.\allowbreak{}cs (\allowbreak{}core package,\allowbreak{} shortened)]{csharp}
public static IReadOnlyProperty? FindTenantPropertyCandidate(
    this IReadOnlyEntityType entityType)
{
    var marked = entityType.GetProperties()
        .Where(p => p[TenancyAnnotationNames.IsTenantProperty] is true)
        .ToList();

    if (marked.Count > 1)
        throw new InvalidOperationException("... several tenant properties ...");
    if (marked.Count == 1)
        return marked[0];

    return typeof(ITenantOwned).IsAssignableFrom(entityType.ClrType)
        ? entityType.FindProperty(nameof(ITenantOwned.TenantId))
        : null;
}
\end{sourcecode}

\begin{datatable}{H{0.92}H{1.10}H{0.98}}
  \textbf{Builder} & \textbf{Typical plugin methods} & \textbf{Note} \\
  \tablehead
  \code{Property\allowbreak{}Builder} / \code{<\allowbreak{}T>} & Column options, roles of
  properties & Here: \code{Is\allowbreak{}Tenant\allowbreak{}Property} \\
  \tablerule
  \code{Index\allowbreak{}Builder} / \code{<\allowbreak{}T>} & Index methods, operator
  classes, filters & Annotation on \code{Metadata} (IMutableIndex) \\
  \tablerule
  \code{Key\allowbreak{}Builder} & Key options (clustering) & As above \\
  \tablerule
  \code{Complex\allowbreak{}Type\allowbreak{}Property\allowbreak{}Builder} & Properties in
  complex types & Only offer where the option makes sense there
  (chapter~\ref{ch:complex-types-and-json}) \\
  \tablerule
  \code{Model\allowbreak{}Builder} & Model-wide options & Annotation on the model \\
\end{datatable}

\section{Runtime annotations}

The SaveChanges interceptor needs to know at run time which entities are isolated and what
their tenant property is called. The convention writes this into an annotation on the entity
type; the interceptor reads it from the \emph{runtime model}:

\begin{sourcecode}[Convention and core package (\allowbreak{}excerpt)]{csharp}
// Convention (design-time model)
entityType.Builder.HasAnnotation(TenancyAnnotationNames.TenantProperty, property.Name);

// Interceptor (runtime model)
public static string? GetTenantPropertyName(this IReadOnlyEntityType entityType)
    => entityType.GetRootType()[TenancyAnnotationNames.TenantProperty] as string;
\end{sourcecode}

Two questions decide whether this works. Does the annotation reach the runtime model? Yes:
\code{Runtime\allowbreak{}Model\allowbreak{}Convention} only removes EF's own design-time
annotations and keeps all others. Does it reach the compiled model? Yes: EF writes simple
values automatically (chapter~\ref{ch:compiled-models-trimming-and-nativeaot}). The custom
\code{ICSharp\allowbreak{}Runtime\allowbreak{}Annotation\allowbreak{}Code\allowbreak{}Generator}
therefore has the opposite job: it removes the RLS annotations that only migrations need:

\begin{sourcecode}[Design/\allowbreak{}Internal/\allowbreak{}Tenancy\allowbreak{}CSharp\allowbreak{}Runtime\allowbreak{}Annotation\allowbreak{}Code\allowbreak{}Generator.\allowbreak{}cs (\allowbreak{}shortened)]{csharp}
public class TenancyCSharpRuntimeAnnotationCodeGenerator
    : NpgsqlCSharpRuntimeAnnotationCodeGenerator
{
    // Constructor: (CSharpRuntimeAnnotationCodeGeneratorDependencies,
    //               RelationalCSharpRuntimeAnnotationCodeGeneratorDependencies) : base(...)

    public override void Generate(IEntityType entityType,
        CSharpRuntimeAnnotationCodeGeneratorParameters parameters)
    {
        if (!parameters.IsRuntime)
        {
            foreach (var key in parameters.Annotations.Keys
                         .Where(k => k.StartsWith(RlsAnnotationNames.Prefix,
                                                  StringComparison.Ordinal))
                         .ToList())
            {
                parameters.Annotations.Remove(key);
            }
        }

        base.Generate(entityType, parameters);
    }
}
\end{sourcecode}

When the compiled model is generated for NativeAOT
(\code{For\allowbreak{}Native\allowbreak{}Aot}), EF also writes the relational model. The
annotation provider has copied the RLS annotations onto the tables, so they appear a second
time there. The generator therefore applies the same filter in an override of
\code{Generate(\allowbreak{}ITable,\allowbreak{}
.\allowbreak{}.\allowbreak{}.\allowbreak{})}.

\begin{datatable}{H{1.36}H{0.88}H{0.88}H{0.88}}
  \textbf{Annotation} & \textbf{Design time} & \textbf{Run time} & \textbf{Compiled model} \\
  \tablehead
  \code{Acme:Rls:*} (policies) & yes (migrations) & yes, unused & filtered \\
  \tablerule
  \code{Acme:Tenancy:Tenant\allowbreak{}Property} & yes & yes (interceptor) & yes \\
  \tablerule
  \code{Acme:Tenancy:Is\allowbreak{}Tenant\allowbreak{}Property} & yes (convention) & yes,
  unused & yes (small) \\
\end{datatable}

\begin{note}[Two kinds of runtime annotations]
  In this book, “runtime annotation” means an ordinary annotation that is read at run time.
  EF also has a runtime annotation API in a narrower sense
  (\code{Add\allowbreak{}Runtime\allowbreak{}Annotation},
  \code{Get\allowbreak{}Or\allowbreak{}Add\allowbreak{}Runtime\allowbreak{}Annotation\allowbreak{}Value}
  on \code{IAnnotatable}). It stores values on the finished model only, for example derived
  data that should not be recomputed on every access. Such values never appear in the
  snapshot or in migrations. For configuration that a plugin sets while the model is built,
  the ordinary annotation is the right choice.
\end{note}

\begin{tip}
  For every annotation it should be clear who reads it and when. Annotations read only at
  design time do not belong in the compiled model; those read at run time must be in it and
  need a test against the generated model. The sample's compiled model test checks both.
\end{tip}

% ------------------------------------------------------------
\chapter{Scaffolding in practice}\label{ch:scaffolding-in-practice}

When scaffolding (\code{dotnet ef dbcontext scaffold}), EF reads the schema and generates
entities and a context. A plugin managing its own schema objects should recognize them in
the process. The route has three stations, two of which the plugin already knows:

\begin{flowdiagram}
  \flowstep{IDatabaseModelFactory}{Reads the schema into a DatabaseModel; the plugin
  attaches RLS data as annotations to the tables}
  \flowstep{Scaffolding model factory}{EF transfers table annotations to the entity types}
  \flowstep{IAnnotationCodeGenerator}{Writes them as HasRowLevelSecurity/HasPolicy
  (chapter~\ref{ch:design-time-services})}
\end{flowdiagram}

\begin{sourcecode}[Scaffolding/\allowbreak{}Tenancy\allowbreak{}Database\allowbreak{}Model\allowbreak{}Factory.\allowbreak{}cs (\allowbreak{}shortened)]{csharp}
public class TenancyDatabaseModelFactory : NpgsqlDatabaseModelFactory
{
    public TenancyDatabaseModelFactory(
        IDiagnosticsLogger<DbLoggerCategory.Scaffolding> logger)
        : base(logger)
    {
    }

    public override DatabaseModel Create(
        DbConnection dbConnection, DatabaseModelFactoryOptions options)
    {
        var model = base.Create(dbConnection, options);
        // ... open the connection if needed ...
        ReadRowLevelSecurity(dbConnection, model);    // pg_class.relrowsecurity / ...force...
        ReadPolicies(dbConnection, model);            // pg_policies
        return model;
    }

    protected virtual void ReadPolicies(DbConnection connection, DatabaseModel model)
    {
        // SELECT schemaname, tablename, policyname, permissive, roles::text[], cmd,
        //        qual, with_check FROM pg_catalog.pg_policies
        // ...
        var policy = new RlsPolicy(
            Name: reader.GetString(2),
            Command: Enum.Parse<PolicyCommand>(reader.GetString(5), ignoreCase: true),
            Permissive: reader.GetString(3) == "PERMISSIVE",
            Roles: roles is ["public"] ? null : roles,    // "public" = default without TO
            Using: reader.IsDBNull(6) ? null : reader.GetString(6),
            WithCheck: reader.IsDBNull(7) ? null : reader.GetString(7));

        table[RlsAnnotationNames.PolicyPrefix + policy.Name] = policy.Serialize();
    }
}
\end{sourcecode}

The factory is registered like the code generator in the design-time services with
\code{Remove\allowbreak{}All}, but with \code{Add\allowbreak{}Scoped}: EF and Npgsql
register the database model factory as scoped, and it depends on the scaffolding logger,
which is scoped as well. The data format is the same as with the Fluent API; snapshot, migrations and code generator therefore work without
further changes.

\begin{note}[Normalized expressions]
  PostgreSQL stores policy expressions in normalized form, for example
  \code{(\allowbreak{}(\allowbreak{}\textquotedbl{}Tenant\allowbreak{}Id\textquotedbl{} =
  (\allowbreak{}NULLIF(\allowbreak{}current\_\allowbreak{}setting(\allowbreak{}'app.\allowbreak{}tenant\_\allowbreak{}id'::text,\allowbreak{}
  true),\allowbreak{} ''::text))::uuid))}. Scaffolded policies contain exactly this text.
  That is correct, but looks different from the form generated by the convention. Teams
  switching to migrations after scaffolding should unify the policies once, otherwise the
  first migration produces a DROP/CREATE with identical content.
\end{note}

\begin{sourcecode}[tests/\allowbreak{}.\allowbreak{}.\allowbreak{}.\allowbreak{}Integration\allowbreak{}Tests/\allowbreak{}Advanced\allowbreak{}Tests.\allowbreak{}cs (\allowbreak{}excerpt)]{csharp}
[Fact]
public void Scaffolding_reads_rls_and_policies()
{
    using var ctx = ContextFor(null);
    var factory = new TenancyDatabaseModelFactory(
        ctx.GetService<IDiagnosticsLogger<DbLoggerCategory.Scaffolding>>());

    var model = factory.Create(
        _fixture.OwnerConnectionString, new DatabaseModelFactoryOptions());
    var orders = model.Tables.Single(t => t.Name == "Orders");

    Assert.Equal(true, orders[RlsAnnotationNames.Enabled]);
    Assert.Equal(true, orders[RlsAnnotationNames.Forced]);
}
\end{sourcecode}

% ------------------------------------------------------------
\chapter{What belongs in the core}\label{ch:what-belongs-in-the-core}

A second provider is the acid test for the package architecture. It shows which parts are
truly provider-independent. The rule of thumb: \textbf{whatever knows no SQL syntax and no
provider type belongs in the core.} For the sample this yields:

\begin{datatable}{H{1.21}H{0.45}H{1.34}}
  \textbf{Building block} & \textbf{Location} & \textbf{Reason} \\
  \tablehead
  \code{Tenant\allowbreak{}Id}, converter, JSON reader/writer & Core & Pure .NET types \\
  \tablerule
  \code{ITenant\allowbreak{}Owned}, \code{ITenant\allowbreak{}Provider}, attributes & Core &
  For entity assemblies without a provider dependency \\
  \tablerule
  \code{Is\allowbreak{}Tenant\allowbreak{}Property(\allowbreak{})}, candidate lookup & Core
  & Metadata only \\
  \tablerule
  Runtime annotation, \code{Get\allowbreak{}Tenant\allowbreak{}Property\allowbreak{}Name} &
  Core & Read by the provider-neutral interceptor \\
  \tablerule
  SaveChanges interceptor & Core & Works only with the change tracker \\
  \tablerule
  \code{Tenant\allowbreak{}Session\allowbreak{}Interceptor} (abstract) & Core & Same flow,
  only the SQL command differs \\
  \tablerule
  Evaluatable filter with attribute & Core & Recognizes function classes of all packages \\
  \tablerule
  Check while loading (materialization interceptor) & Core & Only reads the model, the
  context and the row's values \\
  \tablerule
  Value generation, logging, \code{Add\allowbreak{}Acme\allowbreak{}Tenancy} & Core & No SQL
  involved \\
  \tablerule
  \code{Database.\allowbreak{}Apply\allowbreak{}Tenant\allowbreak{}Session(\allowbreak{})} &
  Core & Uses the provider's interceptor \\
  \tablerule
  Type mapping, translators, migrations, RLS conventions & Provider & Store types, SQL
  syntax, schema objects \\
  \tablerule
  Custom SQL nodes, aggregates & Provider & Depend on syntax and provider factory \\
\end{datatable}

\begin{important}[Name collisions between provider packages]
  Two provider packages that both define an extension method
  \code{Use\allowbreak{}Acme\allowbreak{}Tenancy} on
  \code{Db\allowbreak{}Context\allowbreak{}Options\allowbreak{}Builder} cannot be referenced
  at the same time: the call would be ambiguous. The PostgreSQL package uses EF's builder,
  the SQL Server package the provider's builder
  (\code{Use\allowbreak{}Sql\allowbreak{}Server(\allowbreak{}cs,\allowbreak{} x =>
  x.\allowbreak{}Use\allowbreak{}Acme\allowbreak{}Tenancy(\allowbreak{}))}). If you plan
  several provider packages, use the provider builder consistently from the start, as the
  NetTopologySuite plugin does.
\end{important}

% ------------------------------------------------------------
\chapter{A second provider: SQL Server}\label{ch:a-second-provider-sql-server}

The package
\code{Acme.\allowbreak{}Entity\allowbreak{}Framework\allowbreak{}Core.\allowbreak{}Tenancy.\allowbreak{}Sql\allowbreak{}Server}
implements the same idea with SQL Server's means. The differences show what a plugin really
has to solve anew per database:

\begin{datatable}{H{0.81}H{1.00}H{1.19}}
  \textbf{Aspect} & \textbf{PostgreSQL} & \textbf{SQL Server} \\
  \tablehead
  Tenant in the session & \code{set\_\allowbreak{}config(\allowbreak{}name,\allowbreak{}
  value,\allowbreak{} false)} &
  \code{sp\_\allowbreak{}set\_\allowbreak{}session\_\allowbreak{}context} \\
  \tablerule
  \code{Current\allowbreak{}Tenant(\allowbreak{})} in SQL &
  \code{nullif(\allowbreak{}current\_\allowbreak{}setting(\allowbreak{}.\allowbreak{}.\allowbreak{}.\allowbreak{}),\allowbreak{}
  '')::uuid} & \code{CAST(\allowbreak{}SESSION\_\allowbreak{}CONTEXT(\allowbreak{}N'…') AS
  uniqueidentifier)} \\
  \tablerule
  Column type & \code{uuid} & \code{uniqueidentifier} \\
  \tablerule
  RLS mechanism & \code{CREATE POLICY} per table & Predicate function + \code{CREATE
  SECURITY POLICY} \\
  \tablerule
  Read / write & \code{USING} / \code{WITH CHECK} & \code{FILTER} / \code{BLOCK
  .\allowbreak{}.\allowbreak{}.\allowbreak{} AFTER INSERT/\allowbreak{}UPDATE} \\
  \tablerule
  Table owner & Only bound with \code{FORCE} & No FORCE; predicates are part of the query \\
  \tablerule
  \code{DROP TABLE}, renaming & Policies disappear with it or follow & Policy must go first and be recreated afterwards (SCHEMABINDING) \\
  \tablerule
  DDL in batches & Transactional, freely combinable & \code{CREATE
  FUNCTION/\allowbreak{}SCHEMA} alone in a batch:
  \code{EXEC(\allowbreak{}N'.\allowbreak{}.\allowbreak{}.\allowbreak{}')} \\
  \tablerule
  Table mappings & \multicolumn{2}{>{\RaggedRight\arraybackslash}p{0.6\linewidth}}{Same
  contract from the core package (chapter~\ref{ch:conventions}, page~\pageref{sec:supported-table-mappings})} \\
  \tablerule
  Deliberate exception & Attribute or
  \code{Has\allowbreak{}No\allowbreak{}Row\allowbreak{}Level\allowbreak{}Security(\allowbreak{})}
  & Attribute only \\
\end{datatable}

\section{Activation through the provider builder}

\begin{sourcecode}[Usage and Infrastructure/\allowbreak{}Sql\allowbreak{}Server\allowbreak{}Tenancy\allowbreak{}Options\allowbreak{}Builder\allowbreak{}Extensions.\allowbreak{}cs (\allowbreak{}shortened)]{csharp}
services.AddAcmeTenancy<HttpTenantProvider>();                    // core package
services.AddDbContext<ShopContext>((sp, o) => o
    .UseSqlServer(connectionString, sql => sql.UseAcmeTenancy(sp)));

public static SqlServerDbContextOptionsBuilder UseAcmeTenancy(
    this SqlServerDbContextOptionsBuilder sqlServerOptionsBuilder,
    string sessionKey = "TenantId",
    string policySchema = "acme")
{
    // From the provider builder to the actual DbContextOptionsBuilder
    var optionsBuilder =
        ((IRelationalDbContextOptionsBuilderInfrastructure)sqlServerOptionsBuilder)
        .OptionsBuilder;

    var extension = optionsBuilder.Options.FindExtension<SqlServerTenancyOptionsExtension>()
                    ?? new SqlServerTenancyOptionsExtension();
    // ... validate and set settings ...
    ((IDbContextOptionsBuilderInfrastructure)optionsBuilder)
        .AddOrUpdateExtension(extension.WithSettings(settings));

    return sqlServerOptionsBuilder;
}
\end{sourcecode}

\section{Migrations: security policy instead of policy expression}

For SQL Server, the convention only records the tenant column; the SQL is produced in the
generator. The predicate function is the same for all tables and is created idempotently:

\begin{sourcecode}[Migrations/\allowbreak{}Internal/\allowbreak{}Sql\allowbreak{}Server\allowbreak{}Tenancy\allowbreak{}Migrations\allowbreak{}Sql\allowbreak{}Generator.\allowbreak{}cs (\allowbreak{}excerpt)]{csharp}
protected virtual void CreatePolicy(string table, string? schema, string column,
    MigrationCommandListBuilder builder)
{
    var function = Delimit(PredicateFunctionName, Settings.PolicySchema);
    var target = Delimit(table, schema);
    var predicate = $"{function}({Delimit(column)})";

    // CREATE SCHEMA/FUNCTION must be alone in a batch -> EXEC(N'...')
    Statement(builder, $"IF SCHEMA_ID({Literal(Settings.PolicySchema)}) IS NULL "
        + $"EXEC({Literal($"CREATE SCHEMA {Delimit(Settings.PolicySchema)}")})");

    var functionSql =
        $"CREATE FUNCTION {function}(@TenantId uniqueidentifier) "
        + "RETURNS TABLE WITH SCHEMABINDING AS RETURN SELECT 1 AS [Allowed] "
        + $"WHERE @TenantId = CAST(SESSION_CONTEXT({Literal(Settings.SessionKey)}) "
        + "AS uniqueidentifier)";
    Statement(builder, $"IF OBJECT_ID({Literal(function)}, N'IF') IS NULL "
        + $"EXEC({Literal(functionSql)})");

    Statement(builder,
        $"CREATE SECURITY POLICY {PolicyName(table, schema ?? "dbo")} "
        + $"ADD FILTER PREDICATE {predicate} ON {target}, "
        + $"ADD BLOCK PREDICATE {predicate} ON {target} AFTER INSERT, "
        + $"ADD BLOCK PREDICATE {predicate} ON {target} AFTER UPDATE "
        + "WITH (STATE = ON)");
}

protected override void Generate(DropTableOperation operation, IModel? model,
    MigrationCommandListBuilder builder, bool terminate = true)
{
    // SCHEMABINDING: the policy first, found by the table it binds
    DropPolicy(operation.Name, TableSchema(operation.Schema, model), builder);
    base.Generate(operation, model, builder, terminate);
}
\end{sourcecode}

\code{SCHEMABINDING} only accepts two-part names. A table without an explicit schema is
therefore qualified by the callers with the model's default schema or, failing that, with
\code{dbo}; otherwise SQL Server rejects the security policy.

\begin{important}[Nested SQL]
  The \code{EXEC(\allowbreak{}N'.\allowbreak{}.\allowbreak{}.\allowbreak{}')} contains SQL
  in a string literal that in turn contains literals. The nesting is produced exclusively
  through the string mapping's \code{Generate\allowbreak{}Sql\allowbreak{}Literal}, which
  doubles quotes correctly; identifiers go through \code{Delimit\allowbreak{}Identifier}.
  Session key and schema are additionally validated during configuration.
\end{important}

\section{Schema changes on protected tables}\label{sec:sqlserver-schema-bound-changes}

\code{SCHEMABINDING} protects the predicate function, but it also binds the table. SQL Server
rejects \code{sp\_\allowbreak{}rename} for the table and the tenant column, \code{ALTER SCHEMA
.\allowbreak{}.\allowbreak{}.\allowbreak{} TRANSFER}, \code{ALTER COLUMN} and \code{DROP COLUMN}
as long as a security policy binds the table: “object participates in enforced dependencies”.
This cannot be solved per operation. EF orders a migration's operations itself, and when the
tenant column is renamed, the \code{Rename\allowbreak{}Column\allowbreak{}Operation} comes
\emph{before} the \code{Alter\allowbreak{}Table\allowbreak{}Operation} that carries the changed
annotation. A plain table rename produces no \code{Alter\allowbreak{}Table\allowbreak{}Operation}
at all.

The generator therefore overrides the entry point that receives the whole operation list. It
finds all tables whose bound objects change, follows their renames and wraps the list. At the
front it removes the policies of the affected tables under their old names. At the end it
recreates them with the final table, schema and column names from the target model:

\begin{sourcecode}[Generated SQL when renaming the tenant column (\allowbreak{}shortened)]{sql}
DECLARE @policy nvarchar(400) = (SELECT TOP (1) ... FROM sys.security_policies AS p
    WHERE EXISTS (SELECT 1 FROM sys.security_predicates AS s WHERE s.object_id = p.object_id
        AND s.target_object_id = OBJECT_ID(N'[dbo].[Invoice]')
        AND s.predicate_definition LIKE N'%[[]fn[_]tenant[_]predicate]%'));
IF @policy IS NOT NULL EXEC(N'DROP SECURITY POLICY ' + @policy);
EXEC sp_rename N'[Invoice].[TenantId]', N'Mandant', 'COLUMN';
EXEC sp_rename N'[Invoice].[IX_Invoice_TenantId]', N'IX_Invoice_Mandant', 'INDEX';
IF SCHEMA_ID(N'acme') IS NULL EXEC(N'CREATE SCHEMA [acme]');
IF OBJECT_ID(N'[acme].[fn_tenant_predicate]', N'IF') IS NULL EXEC(N'CREATE FUNCTION ...');
CREATE SECURITY POLICY [acme].[tenant_dbo_Invoice]
    ADD FILTER PREDICATE [acme].[fn_tenant_predicate]([Mandant]) ON [dbo].[Invoice], ...
\end{sourcecode}

Three decisions are built into this:

\begin{itemize}
  \item \textbf{A transaction instead of switching off.} EF runs a migration's commands in a
    transaction, and DDL is transactional in SQL Server. If a step fails, SQL Server also
    rolls back the removal of the policy. A permanently removed policy, or one switched off
    with \code{STATE = OFF}, cannot come about this way. The Down path is the same wrapping
    with the names swapped.
  \item \textbf{Removal by binding, not by name.} Only a policy whose predicate uses the
    plugin's predicate function and that binds exactly this table is dropped. Policies of
    others and tables without RLS stay untouched, and policies under an earlier name are
    found all the same.
  \item \textbf{Schema-qualified names.} Up to 10.0 the policy was called
    \code{tenant\_\allowbreak{}Table}, so tables with the same name in two schemas collided.
    New policies are called \code{tenant\_\allowbreak{}Schema\_\allowbreak{}Table}. Existing
    installations keep the old name until the table is touched the next time; then the
    generator replaces the policy. The remaining ambiguity (schema \code{a\_b} with table
    \code{c} versus schema \code{a} with table \code{b\_c}) and names longer than 128
    characters are rejected by the convention while the model is built.
\end{itemize}

Limits: restoring needs the target model. That is the normal case, since
\code{dotnet ef migrations add} creates a designer file for every migration. If it is
missing, or a changed table does not appear in it, the generator stops instead of dropping a
policy. The application's own schema-bound objects (views with \code{SCHEMABINDING}, its own
security policies) are not handled by the plugin; they still block the same operations. Data
operations that EF sorts between the schema changes run inside the transaction without a
policy, that is with the rights of the migration role.
\code{Sql\allowbreak{}Server\allowbreak{}Migration\allowbreak{}Tests} checks with diffs
generated by EF against a real SQL Server: create, rename column and table, move schema, every
step back and drop, with read and write isolation through the app login after every step. In
addition: tables with the same name in two schemas, a policy under the old name and someone
else's policy on a table without RLS.

\section{Run time: same flow, different command}

\begin{sourcecode}[Runtime/\allowbreak{}Sql\allowbreak{}Server\allowbreak{}Tenant\allowbreak{}Session\allowbreak{}Interceptor.\allowbreak{}cs]{csharp}
public class SqlServerTenantSessionInterceptor : TenantSessionInterceptor
{
    public SqlServerTenantSessionInterceptor(ITenantProvider tenants)
        : base(tenants)
    {
    }

    protected override DbCommand CreateSessionCommand(
        DbConnection connection, DbContext? context, TenantId? tenant)
    {
        var key = context?.GetService<ISqlServerTenancySingletonOptions>()
                      .Settings.SessionKey
                  ?? SqlServerTenancySettings.Default.SessionKey;

        var command = connection.CreateCommand();
        command.CommandText = "EXEC sp_set_session_context @key = @key, @value = @value";
        AddParameter(command, "@key", key);
        AddParameter(command, "@value", tenant?.Value);    // NULL: no rows visible
        return command;
    }
}
\end{sourcecode}

\section{Deliberately not ported}

The SQL Server package is a complete but leaner provider. Fluent API policies, the optional
query filter, custom SQL nodes, aggregates, the snapshot code generator and scaffolding are
missing there. That is a product decision, not a technical obstacle: each of these chapters
can be transferred to SQL Server with the same patterns. The list belongs in the package's
documentation so users know what to expect.

\begin{sourcecode}[tests/\allowbreak{}.\allowbreak{}.\allowbreak{}.\allowbreak{}Sql\allowbreak{}Server.\allowbreak{}Tests/\allowbreak{}Sql\allowbreak{}Server\allowbreak{}Tests.\allowbreak{}cs (\allowbreak{}excerpt)]{csharp}
[Fact]
public async Task Block_predicate_rejects_rows_of_other_tenant()
{
    await using var db = ContextFor(TenantId.New());

    var ex = await Assert.ThrowsAsync<SqlException>(() => db.Database.ExecuteSqlAsync(
        $"INSERT INTO [Orders] ([TenantId], [Total]) VALUES ({Guid.NewGuid()}, 1)"));

    Assert.Equal(33504, ex.Number);      // violation of a BLOCK PREDICATE
}
\end{sourcecode}

As with PostgreSQL, the tests run on two levels: without a database for model, query and
migration SQL, and with Testcontainers against a real SQL Server. There, the application
connects with its own login instead of \code{sa} to reproduce the production situation.
